\section{Introduction}

Euclidean tangent spaces describe the geometry of a metric space at a
point, but a coordinate chart must compare all pairs of points in a
neighborhood. Under a lower Ricci curvature bound these are distinct
regularity questions.

Throughout, noncollapsed means that the reference measure is exactly
\(\mathcal H_d^n\). We use the full \(\CD(K,n)\) condition together with
quadratic Cheeger energy. Spaces are complete, separable, geodesic, and
have full-support reference measure finite on bounded sets. A point is
\emph{\(n\)-regular} if every pointed Gromov--Hausdorff tangent there is
pointed-isometric to \((\R^n,|\cdot|,0)\). Here a tangent is a pointed
limit of rescalings by inverse radii tending to infinity.
We prove the following open-neighborhood result.

\begin{theorem}\label{thm:main}
For every integer \(n\ge2\) there is a constant \(L_n\ge1\) with
the following property. Let \(K\in\R\), and let
\((X,d,\mathcal H_d^n)\) be a noncollapsed \(\RCD(K,n)\) space.
Let \(p\in X\).
If every pointed Gromov--Hausdorff tangent of \(X\) at \(p\) is
pointed-isometric to \((\R^n,|\cdot|,0)\), then there exist an open
neighborhood \(U\) of \(p\), an open set \(V\subset\R^n\), and a
homeomorphism \(F:U\to V\) such that
\[
 L_n^{-1}d(x,z)\le |F(x)-F(z)|\le L_nd(x,z)
 \qquad (x,z\in U).
\]
The distance on the left and right is the original distance \(d\)
restricted to \(U\).
\end{theorem}

The neighborhood may depend on \(X\) and \(p\); the distortion
constant depends only on the dimension. The theorem concerns every
prescribed regular point, not only almost every point of the reference
measure. It also yields the following manifold region.

\begin{corollary}\label{cor:full-measure-atlas}
Under the hypotheses of Theorem~\ref{thm:main}, there is an open set
\(\mathcal U\subset X\) containing every \(n\)-regular point, with
\(\mathcal H_d^n(X\setminus\mathcal U)=0\), and a countable atlas
of \(L_n\)-bi-Lipschitz charts from ambient-open subsets of
\(\mathcal U\) onto open subsets of \(\R^n\). Each chart uses the
restricted ambient distance, and its transition maps are
\(L_n^2\)-bi-Lipschitz.
\end{corollary}

\begin{proof}
The \(n\)-regular set has full \(\mathcal H_d^n\) measure by
\cite[Theorem~1.12(ii),(vi)]{DPG18}, applied with reference measure
\(\mathcal H_d^n\). Take \(\mathcal U\) to be the union of the chart
domains in Theorem~\ref{thm:main}. As a separable metric space,
\(\mathcal U\) is second countable, so these domains admit a
countable subcover. On an overlap, composing one chart with the
inverse of the other gives the stated transition bound.
\end{proof}

The open region \(\mathcal U\) may contain points outside the exact
regular locus; the corollary does not assert that this locus is open.

\subsection{History and significance}

The structure theory of noncollapsed limits of Riemannian manifolds
with lower Ricci bounds provides an early form of this question.
Cheeger and Colding obtained manifold neighborhoods with bi-H\"older
parametrizations and asked about a bi-Lipschitz upgrade
\cite[Theorem~5.14 and Remark~5.15]{CC97}. Their question concerns
quantitatively regular Ricci-limit neighborhoods and is a historical
antecedent of the synthetic regular-point problem considered here.

Synthetic lower Ricci bounds were developed through optimal transport
by Lott and Villani and by Sturm \cite{LV09,Sturm06II}. The
Riemannian refinement of Ambrosio, Gigli and Savar\'e distinguishes
quadratic Cheeger energy and linear heat flow \cite{AGS14}; the
finite-dimensional connections with Bochner inequalities were
developed by Erbar, Kuwada and Sturm \cite{EKS15}. In the noncollapsed
setting, De Philippis and Gigli established volume convergence,
volume rigidity, and quantitative Euclidean approximation
\cite{DPG18}.

Mondino and Naber proved countable bi-Lipschitz rectifiability of
RCD spaces: measurable pieces cover almost every point and map to
measurable Euclidean subsets \cite[Theorem~1.1]{MN19}. This does not
give an ambient-open neighborhood of each prescribed regular point.
Kapovitch and Mondino proved that the entire regular set is contained
in an open manifold region with bi-H\"older charts
\cite[Theorem~3.1 and Corollary~3.2]{KM21}. This statement does not
require the exact regular locus itself to be open. An open
bi-Lipschitz chart requires stronger control across infinitely many
scales.

The conjecture that every regular point in the noncollapsed setting
admits a bi-Lipschitz chart is explicitly recalled in the recent work
of Honda and Zhang \cite[Section~1.2]{HZ26}. Theorem~\ref{thm:main}
answers this question under the precise noncollapsed \(\RCD(K,n)\)
hypotheses stated above. Honda and Zhang's geometric regularity
theorems assume a positive lower bound for injectivity radius on
neighborhoods \cite[Theorems~1.1 and 3.35]{HZ26}. Their analytic
criterion also permits Morrey \(L^2\) control of the Laplacians,
together with a neighborhood-uniform positive Gram determinant
\cite[Theorem~3.9 and Proposition~3.10]{HZ26}. Neither additional
hypothesis is assumed here.

Several nearby results clarify the strength of the conclusion.
Colding and Naber proved that bounded subsets of quantitative
Reifenberg loci in noncollapsed Ricci limits, with sufficiently small
flatness and scale parameters, admit bi-Lipschitz embeddings into a
Euclidean space of possibly higher dimension
\cite[Theorem~1.1]{CN13}. They also constructed examples with
Euclidean tangents that admit no coordinates inducing a H\"older
continuous metric tensor \cite[Theorem~1.2]{CN13}. The conclusion here
is an open chart into \(\R^n\); it does not assert H\"older continuity
of the metric tensor in that chart. Likewise, the restrictions on
bi-Lipschitz regular maps in \cite[Proposition~5.1]{HP23} concern
maps whose pullbacks of smooth target test functions have locally
bounded Laplacians \cite[Definition~3.1]{HP23}. No such analytic
regularity is asserted for the limiting coordinates constructed
below; in particular, their components are not asserted to have
bounded Laplacian.

\subsection{The two main estimates}

After rescaling at \(p\), volume comparison makes all sufficiently
small balls at nearby centers quantitatively Euclidean. Harmonic
functions approximating Euclidean coordinates then give an open
coordinate map \(u=(u_1,\ldots,u_n)^T\). The scale-dependent mean Gram matrices of \(u\)
are invertible, but they need not be uniformly nondegenerate as the
scale tends to zero. We compensate for this degeneration in the
\emph{space of coordinate values}.

Two kinds of control make this possible.  Accumulated Bochner mass
records how much the coordinates must be stretched to compensate for
their degeneration; this accumulation may be unbounded.  A separate
estimate evaluates the same mass forms on actual coordinate differences
$u(x)-u(z)$ and controls them by defects whose sum is small.  The
construction uses the cumulative mass to prescribe the stretching and
the directional estimate to keep its metric error summable.

Write \(r_i=2^{-i}\), and let
\(G=(\langle\nabla u_j,\nabla u_k\rangle)_{j,k=1}^n\) be the Gram
matrix field. Its mean on \(B_{r_i}(y)\) is denoted by \(R_i(y)\).
Let \(\mathfrak m\) denote Hausdorff measure in the current length
units. If the rescaled curvature lower bound is \(-\kappa\), the
matrix-valued Bochner measure is
\(M=\tfrac12\Delta G+\kappa G\mathfrak m\), with the Laplacian
understood entrywise. For a fixed enlargement \(L>1\), define the
positive semidefinite matrix
\[
 D_{i,y}=R_i(y)^{-1}
       \bigl[r_i^{2-n}M(B_{Lr_i}(y))\bigr]R_i(y)^{-1},
 \qquad D[\xi]=\xi^TD\xi.
\]
Corrected heat averages of \(G\) yield positive definite forms
\(H_i(x)\), increasing with \(i\), on the space of coordinate
values. Lemma~\ref{co:cumulative-comparison} compares these forms
with accumulated bounded-overlap sums of the local matrices
\(D_{i,y}\), including a reverse comparison for suitable weighted
sums. This is the cumulative input to the construction.

The first main estimate gives directional control of the same local
matrices. For each fixed \(L,C_{\mathrm{sep}}>1\), there are
nonnegative defect functions \(E_i\) on the working region such that,
whenever \(x,z,y\in B_4(p)\) and
\(d(x,y),d(z,y)\le C_{\mathrm{sep}}r_i\),
\[
 D_{i,y}[u(x)-u(z)]
 \le C r_i^2\bigl(E_i(x)+E_i(z)\bigr),
 \qquad \sup_{x\in B_4(p)}\sum_{i\ge0}E_i(x)\longrightarrow0
\]
as the initial rescaling becomes flat
(Proposition~\ref{pi:directional-pinching}). The proof uses Poisson
replacements of squared distance, affine approximation of radial
derivatives, and a sharp projection estimate for the difference of
two replacements. A radial transport argument preserves the
\emph{squared} defect in this last estimate. Poisson replacement,
harmonic transformation, heat Gram normalization, and sharp
squared-Hessian pinching have substantive antecedents in the work
of Cheeger and Colding and of Cheeger, Jiang and Naber
\cite{CC96,CJN21}. The two-center directional bound displayed above
is proved here in the form needed for the value-space construction.

Scale-dependent transformations are also central in
\cite[Theorem~1.32]{CN15} and \cite[Theorems~1.7 and 4.1]{HH24}.
Potentially unbounded, noncommuting normalization matrices already
occur in these methods and in \cite[Theorem~3.8]{HZ26}; those
freedoms alone do not remove the H\"older loss. The issue here is
how to use the directional estimate to control the cumulative
metric error without assuming summability of its square root.

The second main ingredient is the analytic construction in
Proposition~\ref{re:metric-iteration}. On the common open value
domain \(\Omega=u(B_1(p))\), the construction produces smooth maps
\(h_i:\Omega\to\R^n\) whose ordinary Jacobians \(Dh_i\) remain
invertible. From the forms \(D_{i,y}\) and smooth cutoffs we build
covector fields \(v_i\) and, for \(P\in\Omega\), iterate
\[
 h_0(P)=P,\qquad
 h_{i+1}(P)=h_i(P)+Dh_i(P)^{-T}v_i(P).
\]
The leading change in the pullback metric is a positive sum of the
forms \(D_{i,y}\). When the inverse transpose is differentiated, its
linear contribution to a squared length contains
\(D^2h_i(\xi,\xi)\). A recurrence for these diagonal derivatives
allows the errors to be summed using \(\sum_i E_i\), even when
\(\sum_i\sqrt{E_i}\) diverges. A separate estimate for all derivatives,
with Gevrey growth in their order, closes the iteration despite the
loss of one derivative in each update.

This part of the argument uses only positive matrix comparisons,
directional smallness, and smooth cutoffs; it does not differentiate
the scale metrics \(H_i(x)\) or assume they remain bounded. It isolates
an analytic way to turn summable directional information into
coordinates, separate from the geometric origin of that information.
The limit \(F=\lim_i h_i\circ u\) has the asserted metric bounds.
The openness of its image follows from the open coordinates \(u\)
and invariance of domain.

Summability also governs bi-Lipschitz Reifenberg parametrizations.
For extrinsic sets, David and Toro control coherent changes of
approximating planes by square-summable errors
\cite[Proposition~8.34]{DT12}. For metric spaces, Gigli and Violo
obtain bi-Lipschitz parametrizations under summability of ordinary
dyadic Gromov--Hausdorff errors \cite[Theorems~2.2 and 2.7]{GV25}.
Regularity at a point does not itself supply that summability.
Here the summable quantity is instead the directional Bochner
defect, and its conversion into metric control is proved by the
explicit iteration above.

Section~\ref{co:section} constructs the open harmonic domain and
its scale normalizations. Section~\ref{co:heat-section} proves the
cumulative matrix comparison, and Section~\ref{pi:section} proves
the summable directional estimate. Section~\ref{re:section} uses
these two inputs to construct the renormalization and prove the
all-pairs metric bounds and openness of the limiting chart.

\subsection{Conventions}

All balls use the ambient metric. Every enlargement appearing in a
local estimate is fixed before a sufficiently late rescaling is
chosen. We write \(C\) for constants depending on the dimension and
on specified fixed enlargement factors. Uniform quantities tending
to zero along the rescalings will be denoted by \(o(1)\); this
notation includes uniformity over the centers and all dyadic scales
under consideration. A subscript, as in \(o_A(1)\), permits
dependence on constants fixed earlier.

For a positive definite matrix \(A\), set
\[
 A[\xi]=\xi^T A\xi,\qquad |\xi|_A=A[\xi]^{1/2}.
\]
Matrix inequalities are inequalities of quadratic forms. Vectors
and covectors are represented using the fixed Euclidean pairing,
and their length norms are specified explicitly. The norm
\(\|T\|_{H\to Q}\) of a multilinear map uses \(|\cdot|_H\) on each
argument and \(|\cdot|_Q\) on its value. Derivatives of maps on
coordinate value space are ordinary Euclidean derivatives. We use
\(\fint_B\) for the integral divided by the measure of \(B\).
